% =====================================================================
% P1 Appendix A.5 --- Proposition 2 (branch-wise calibration and its price).
%
% Source of truth: drafts/proofs-props8-12.md v1.4, Sections 5, 7, 8
% (FS-10, FS-12, FS-12.1).  This file is a fragment: it is meant to be
% \input under \appendix \section{Proofs} after A.4.
%
% Theorem-like environments use the explicit paper numbering
% (Propositions 2A--2C).  Equation tags (5.1)--(5.3), (7.1)--(7.2)
% follow the frozen source; do not renumber without updating the main
% text.
% =====================================================================

\newtheorem*{aVpropA}{Proposition 2A}
\newtheorem*{aVpropB}{Proposition 2B}
\newtheorem*{aVpropC}{Proposition 2C}

% Fixed equation tags: keep the frozen printed numbers while giving
% hyperref a unique destination for every display.
\providecommand{\theHequation}{}
\let\aVtheHequation\theHequation
\newcounter{aVanchor}
\renewcommand{\theHequation}{aV.\arabic{aVanchor}}
\newcommand{\aVfixedtag}[1]{\stepcounter{aVanchor}\tag{#1}}

\subsection{Branch-wise calibration and its price (Proposition~2)}\label{app:a5}

This appendix gives the finite-sample statements and proofs behind
Proposition~2.  Proposition~2A is the branch-wise split-conformal
coverage guarantee, Proposition~2B the coverage dominance of the
branch-union threshold, and Proposition~2C the absence of a
distribution-free coverage gain per unit width.  The notation of
Appendix~A.4 is in force.

\paragraph{Setting.}
Let $D_n=\{(X_i,Y_i,B_i)\}_{i=1}^n$ be a calibration sample and let
$(X_{n+1},Y_{n+1},B_{n+1})$ be a test point.  The calibration sample and
the test point are exchangeable, and conditional on $B=b$ the pairs
$(X_i,Y_i)$ with $B_i=b$ are exchangeable with the test pair.  Write
\[
n_b=\#\{i\le n:B_i=b\},
\qquad
k_b=\bigl\lceil(n_b+1)(1-\alpha)\bigr\rceil.
\]
If $k_b\le n_b$, let $\hat q_b$ be the $k_b$-th order statistic of the
branch-$b$ calibration scores; if $k_b>n_b$, set $\hat q_b=+\infty$.
Put
\[
C_b(x)=\{y:s(x,y)\le\hat q_b\},
\qquad
C_\cup(x)=C_0(x)\cup C_1(x).
\]
For the dominance statement we also need the marginal threshold: let
$k=\lceil(n+1)(1-\alpha)\rceil$, let $\hat q$ be the $k$-th order
statistic of all $n$ calibration scores (with $\hat q=+\infty$ if
$k>n$), and put $\hat Q=\max(\hat q_0,\hat q_1)$.

\subsubsection{Branch-wise coverage (Proposition~2A)}

\begin{aVpropA}
Under the exchangeability assumption above, for $b\in\{0,1\}$,
\begin{equation}\label{eq:a5-5.1}
\Pr\bigl(Y_{n+1}\in C_b(X_{n+1})\mid B_{n+1}=b\bigr)\ge1-\alpha,
\aVfixedtag{5.1}
\end{equation}
and
\begin{equation}\label{eq:a5-5.2}
\Pr\bigl(Y_{n+1}\in C_\cup(X_{n+1})\mid B_{n+1}=b\bigr)\ge1-\alpha.
\aVfixedtag{5.2}
\end{equation}
If $\hat q_b=+\infty$, then $C_b=\mathcal Y$, the coverage is $1$, and
the inequality still holds.
\end{aVpropA}

\begin{proof}
Condition on the branch labels and on $n_b$.  The $n_b$ calibration
scores in branch $b$ and the test score are $n_b+1$ exchangeable random
variables.  Let $R\in\{1,\dots,n_b+1\}$ be the rank of the test score
among them, with random tie-breaking; $R$ is uniform on
$\{1,\dots,n_b+1\}$.  Since $\hat q_b$ is the $k_b$-th order statistic,
\[
Y_{n+1}\in C_b(X_{n+1})
\iff s(X_{n+1},Y_{n+1})\le\hat q_b,
\]
and $R\le k_b$ implies $s(X_{n+1},Y_{n+1})\le\hat q_b$.  With ties,
$R\le k_b$ still implies $s\le\hat q_b$.  Therefore
\begin{equation}\label{eq:a5-5.3}
\begin{aligned}
\Pr\bigl(Y_{n+1}\in C_b(X_{n+1})\mid B_{n+1}=b,\,n_b\bigr)
&\ge\Pr(R\le k_b)\\
&=\frac{k_b}{n_b+1}
\ge1-\alpha,
\end{aligned}
\aVfixedtag{5.3}
\end{equation}
where the last inequality uses
$k_b=\lceil(n_b+1)(1-\alpha)\rceil\ge(n_b+1)(1-\alpha)$.  Taking
expectations over $n_b$ gives \eqref{eq:a5-5.1}.  Since
$C_b\subseteq C_\cup$, \eqref{eq:a5-5.2} follows immediately from
\eqref{eq:a5-5.1}.
\end{proof}

\paragraph{Conditioning level.}
The guarantee is branch-conditional (on $B=b$) but marginal in
$(X,Y)$: it is not an $X$-conditional guarantee.  Like the
$\kappa_{\mathrm{conf}}$ bound of Theorem~2, it cannot be read as
coverage at each $x$.

\paragraph{Union, convex hull, and width.}
For a symmetric score
$s(x,y)=|y-\mu(x)|/\mathrm{scale}(x)$,
\[
C_b(x)=[\mu(x)-\hat q_b\,\mathrm{scale}(x),\,
\mu(x)+\hat q_b\,\mathrm{scale}(x)],
\]
and $C_\cup(x)$ is the interval with radius
$\hat Q=\max(\hat q_0,\hat q_1)$.  For CQR,
\[
C_b(x)=[\hat q_{\mathrm{lo}}(x)-\hat q_b,\,
\hat q_{\mathrm{hi}}(x)+\hat q_b],
\]
and the convex hull of $C_\cup(x)$ is
$[\hat q_{\mathrm{lo}}(x)-\hat Q,\,
\hat q_{\mathrm{hi}}(x)+\hat Q]$.  Under CQR both branch intervals
contain the base interval
$[\hat q_{\mathrm{lo}}(x),\hat q_{\mathrm{hi}}(x)]$, so the union is
itself an interval and equals its convex hull; for a general score the
union need not be an interval, and the convex hull is the safe object.
For CQR (and for a symmetric score with base $0$),
\[
w_\cup(x)=\mathrm{base}(x)+2\hat Q\,\mathrm{scale}(x)
\ge\mathrm{base}(x)+2\hat q\,\mathrm{scale}(x)=w_{\mathrm{marg}}(x),
\]
where $\hat q$ is the marginal threshold of Proposition~2B.  The extra
width is $2(\hat Q-\hat q)$ in the CQR case, or the corresponding
scale factor in the symmetric case.

\paragraph{Small branches and infinite thresholds.}
If $k_b>n_b$ (for example, $n_b=0$, or $n_b$ small enough that the
$1-\alpha$ quantile lies beyond the sample), then
$\hat q_b=+\infty$, $C_b=\mathcal Y$, the coverage is $1$, and the width
is $+\infty$.  This is a legitimate finite-sample outcome, not an
implementation error.  Accordingly, we report $n_b$ and the number of
cells with $\hat q_b=+\infty$, and exclude infinite widths from average
widths.  In the experiments the smallest branches contain
only a few dozen calibration points per fold, so their width ratios are
dominated by high-quantile noise; they should not be read as
population-level inflation factors.  Section~5.2 reports the branch counts
for the branch-union runs (112 branch cells, smallest $n_b=16$, no
infinite threshold).

\subsubsection{Coverage dominance (Proposition~2B)}

\begin{aVpropB}
Let $s_1,\dots,s_n$ be the calibration scores,
$k=\lceil(n+1)(1-\alpha)\rceil$, and let $\hat q$ be the $k$-th order
statistic (with $\hat q=+\infty$ if $k>n$).  Split the scores by branch
$b\in\{0,1\}$ into $n_b$ scores, set
$k_b=\lceil(n_b+1)(1-\alpha)\rceil$, and let $\hat q_b$ be the
$k_b$-th order statistic of branch $b$ (with $\hat q_b=+\infty$ if
$k_b>n_b$).  Put $\hat Q=\max(\hat q_0,\hat q_1)$.  Then, for any
calibration set, with ties and empty branches allowed,
\begin{equation}\label{eq:a5-7.1}
\hat Q\ge\hat q.
\aVfixedtag{7.1}
\end{equation}
Consequently, for every test point $(x,y)$, $s(x,y)\le\hat q$ implies
$s(x,y)\le\hat Q$; the set $\{y:s(x,y)\le\hat Q\}$ contains the marginal
conformal set $\{y:s(x,y)\le\hat q\}$.  For a common scalar score this
set is exactly $C_0(x)\cup C_1(x)$.  For CQR the union is an interval
and equals its convex hull; for a general score the union may not be an
interval, so coverage and width comparisons are stated for its convex
hull.  If the branches are calibrated with different score functions,
the raw union is not a sublevel set of a common score; dominance of the
marginal set is then not automatic and must be checked separately.
\end{aVpropB}

\begin{proof}
We first prove the ceiling inequality
\begin{equation}\label{eq:a5-7.2}
k=\lceil(n+1)c\rceil
\le\lceil(n_0+1)c\rceil+\lceil(n_1+1)c\rceil
=k_0+k_1,
\qquad c=1-\alpha>0.
\aVfixedtag{7.2}
\end{equation}
Indeed, $\lceil x\rceil+\lceil y\rceil\ge\lceil x+y\rceil$, and
\[
(n_0+1)c+(n_1+1)c=(n+2)c\ge(n+1)c.
\]
Since the ceiling is monotone, this gives \eqref{eq:a5-7.2}.

\emph{Case 1: $\hat q=+\infty$.}  Then $k>n$.  By \eqref{eq:a5-7.2},
$k_0+k_1\ge k>n=n_0+n_1$, so at least one $k_b>n_b$; hence
$\hat q_b=+\infty$ and $\hat Q=+\infty\ge\hat q$.

\emph{Case 2: $\hat q<+\infty$.}  Suppose for contradiction that
$\hat Q<\hat q$.  For $t\in\mathbb R$, let
$N(t)=\#\{i:s_i\le t\}$ and
$N_b(t)=\#\{i:B_i=b,\ s_i\le t\}$.  Since $\hat q$ is the $k$-th order
statistic, every $t<\hat q$ satisfies $N(t)\le k-1$; ties do not affect
this count.  Taking $t=\hat Q<\hat q$ gives $N(\hat Q)\le k-1$.  On the
other hand, for each $b$:
\begin{itemize}
\item if $\hat q_b<+\infty$, then $\hat q_b\le\hat Q$, and the $k_b$-th
order statistic being at most $\hat Q$ means that at least $k_b$
branch-$b$ scores are at most $\hat Q$, that is,
$N_b(\hat Q)\ge k_b$;
\item if $\hat q_b=+\infty$, then $\hat Q=+\infty$, contradicting
$\hat Q<\hat q<+\infty$.
\end{itemize}
Therefore
\[
N(\hat Q)=N_0(\hat Q)+N_1(\hat Q)\ge k_0+k_1\ge k
\]
by \eqref{eq:a5-7.2}, contradicting $N(\hat Q)\le k-1$.  Hence
$\hat Q\ge\hat q$.
\end{proof}

\paragraph{Dominance and width.}
\begin{itemize}
\item For a common score, the marginal set is contained in the union
pointwise, so coverage is monotone.
\item For CQR,
$C_{\mathrm{marg}}(x)=[\hat q_{\mathrm{lo}}(x)-\hat q,\,
\hat q_{\mathrm{hi}}(x)+\hat q]$, while the convex hull of $C_\cup(x)$
is $[\hat q_{\mathrm{lo}}(x)-\hat Q,\,
\hat q_{\mathrm{hi}}(x)+\hat Q]$, which contains
$C_{\mathrm{marg}}(x)$.  Under CQR the two branch intervals both
contain the base interval, so the union is an interval and equals its
convex hull; for a general score the union may not be an interval.  The
convex hull contains the union, so its coverage is at least the union
coverage and at least the coverage of each $C_b$.
\item The width comparison is
$w_\cup(x)=\mathrm{base}(x)+2\hat Q\,\mathrm{scale}(x)
\ge\mathrm{base}(x)+2\hat q\,\mathrm{scale}(x)=w_{\mathrm{marg}}(x)$;
the CQR extra width is $2(\hat Q-\hat q)$.
\item No within-branch exchangeability is required for
\eqref{eq:a5-7.1}: it is a pure order-statistic inequality that holds
for every fixed calibration set.
\end{itemize}

\subsubsection{No distribution-free gain per unit width (Proposition~2C)}

\begin{aVpropC}
Fix a calibration set with $\hat Q>\hat q$.  For every test sample size
$m\ge1$ and every positive function $g(\Delta w)$, there exists a test
distribution such that:
\begin{enumerate}
\item[(i)] the population branch-$b$ coverage gain is
\[
G_b^{\mathrm{pop}}
=F_{\mathrm{test},b}(\hat Q)-F_{\mathrm{test},b}(\hat q)=0,
\]
and, for every $m$, the empirical gain
\[
\hat G_{b,m}
=\hat F_{\mathrm{test},b}(\hat Q)-\hat F_{\mathrm{test},b}(\hat q)=0
\]
almost surely (with the convention that the empirical gain is $0$ when
branch $b$ has no test point);
\item[(ii)] the width increment is
$\Delta w=2(\hat Q-\hat q)\,\mathrm{scale}>0$ for a symmetric score, and
$\Delta w=2(\hat Q-\hat q)$ for CQR.
\end{enumerate}
Hence no distribution-free lower bound of the form
$G_b\ge g(\Delta w)>0$ can hold.  The width cost is computable from the
calibration set; the gain is not.
\end{aVpropC}

\begin{proof}
Take a test distribution whose branch-$b$ score is supported on
$(-\infty,\hat q]$ (for example, a point mass $s_0\le\hat q$ in each
branch, with branch probability $q_b$).  Then
\[
F_{\mathrm{test},b}(\hat Q)
=F_{\mathrm{test},b}(\hat q)=1,
\qquad
G_b^{\mathrm{pop}}=0,
\]
while $\hat Q>\hat q$ gives $\Delta w>0$.  For every $m$, all test
scores are at most $\hat q$, so
$\hat F_{\mathrm{test},b}(\hat Q)
=\hat F_{\mathrm{test},b}(\hat q)=1$ and the empirical gain is $0$
almost surely.  If the test distribution is required to be continuous,
approximate the point mass by a narrow density supported on
$(-\infty,\hat q]$; the gain is still $0$ (or arbitrarily small), while
the width increment is unchanged.  Thus every positive $g(\Delta w)$ is
violated by this test distribution.
\end{proof}

\paragraph{Scope.}
\begin{itemize}
\item Proposition~2C does not claim that branch-union is useless in
practice.  It does restore branch coverage (Proposition~2A); the gain
depends on the target distribution.
\item It does not claim that the width cost is negligible.  The width
increment is computable from the calibration set and must be reported.
\item It does not claim that every test distribution gives zero gain.
The existence of one zero-gain distribution is enough to rule out a
distribution-free lower bound.
\item The result is the width-axis counterpart of the non-identification
in Theorem~3B: both express the same missing information, one through
coverage and one through width.
\end{itemize}

\let\theHequation\aVtheHequation
